\subsection{当 $Y$ 为超 Kähler 且 $M$ 紧时会发生什么}\label{hyperm}

在结束这个主题之前，我们不应忘记说明：在一个重要的特殊情形中，存在一种更简单的办法，可以部分地比较膜量子化与几何量子化。

在整个分析中，我们只是把 $Y$ 当作复辛流形。我们并未假设 $Y$ 的复辛结构能延拓为第 \ref{amo} 节所述意义下的完备超 Kähler 结构。这个假设足以保证 $Y$ 存在良好的 $A$-模型，并且它在许多有趣的例子中成立。但它很可能并非必需。

若我们确实假设 $Y$ 的复辛结构延拓为超 Kähler 结构，那么就会出现一种更简单的办法来部分比较膜量子化与几何量子化。这其实已在 \cite{GW}（第 2.3 与 4 节）中讨论过。首先，$Y$ 的对称 $\tau$ 在复结构 $I$ 中是反全纯的，但它在复结构 $J$ 中实际上是全纯的（见式 \eqref{tabush}）。因此 $M$ 作为不动点集的一个分支，在复结构 $J$ 中是全纯的。由于 $\omega_J$ 关于 $J$ 是 Kähler 的，这意味着复结构 $J$ 给出 $M$ 的一个复 Kähler 极化，即几何量子化意义下的极化。

用这个极化对 $M$ 作几何量子化，得到的 Hilbert 空间是
\be\label{hgq}\H_\gq=H^0(M,K_M^{1/2}\otimes\frL),\ee
其中照例 $\frL$ 是预量子线丛。（为简单起见，我们假设 $\frL$ 的 $\CP$ 曲率 $F=\omega_J$ 足够正，使高阶上同调消失。）Hilbert 空间结构由内积
\be\label{ngg}(\psi,\chi)=\int_M\d\mu\,\bar\psi\chi\ee
给出，其中 $\d\mu$ 是由辛结构 $\omega_J$ 导出的测度。

让我们看看能否为 $\H_\gq$ 与膜量子化所得之间建立关系。首先注意，在此情形中，$\B_\cc$ 与支撑在 $M$ 上的拉格朗日膜 $\B$ 都是 $(A,B,A)$ 型膜：它们关于辛结构 $\omega_I$ 与 $\omega_K$ 是 $A$-膜，关于复结构 $J$ 是 $B$-膜。这可如下看出。\footnote{作为 $\omega_I$、$\omega_K$ 的 $A$-模型以及 $J$ 的 $B$-模型的微分的超荷是线性相关的。因此只需检验其中两个微分下的不变性；第三个不变性随之成立。直接检验全部三个不变性似乎更有启发性，所以我们这样做。}

对 $\B$ 的情形，我们只需注意：由于 $\omega_I$ 是 $\tau$-奇的（式 \eqref{tabush}），$\tau$-不变子流形 $M$ 关于 $\omega_I$ 是拉格朗日的，正如它关于 $\omega_K$ 一样。因此 $\B$ 关于 $\omega_I$ 是 $A$-膜，正如它关于 $\omega_K$ 一样。另一方面，我们已经注意到 $M$ 在复结构 $J$ 中是复子流形。$\B$ 的 $\CP$ 丛满足 $\sF+\sB=0$。特别地，$\sF+\sB$ 在复结构 $J$ 中为 $(1,1)$ 型。支撑在复子流形上、且 $\sF+\sB$ 为 $(1,1)$ 型的膜是 $B$-膜，故 $\B$ 在复结构 $J$ 中是 $B$-膜。简言之，它是 $(A,B,A)$ 型膜。

对 $\B_\cc$，其支撑是整个 $Y$，而 $Y$ 在复结构 $J$ 中是复流形；其 $\CP$ 曲率满足 $\sF+\sB=\omega_J$，后者在复结构 $J$ 中为 $(1,1)$ 型。这些条件同样保证 $\B_\cc$ 关于 $J$ 是 $B$-膜。$\B_\cc$ 关于 $\omega_K$ 是 $A$-膜，因为 $\omega_K^{-1}(\sF+\sB)=I$ 是可积复结构；类似地，它关于 $\omega_I$ 是 $A$-膜，因为 $\omega_I^{-1}(\sF+\sB)=-K$ 是可积复结构。所以 $\B_\cc$ 同样是 $(A,B,A)$ 型膜。

对于 $M$ 的量子化，按我们通常的 $\CP$ 丛选择（$\B_\cc$ 取平凡丛，$\B$ 取 $\frL^{-1}$），膜量子化给出的答案是 $\omega_K$ 的 $A$-模型中 $(\B,\B_\cc)$ 弦的空间。我们一直简单地称之为 $\Hom(\B,\B_\cc)$，但在当前语境中最好写成
$\Hom_{A_K}(\B,\B_\cc)$，下标 $A_K$ 提醒我们考虑的是 $\omega_K$ 型 $A$-模型中的物理弦态。几何量子化用 Kähler 极化 $J$ 对 $M$ 量子化给出的答案 \eqref{hgq}，等同于 $\Hom_{B_J}(\B,\B_\cc)$，也就是复结构 $J$ 的 $B$-模型中 $(\B,\B_\cc)$ 弦的空间。

因此，用复结构 $Y$ 的膜量子化，与用复 Kähler 极化 $K$ 的几何量子化一致，当且仅当
$\Hom_{A_K}(\B,\B_\cc)$ 与 $\Hom_{B_J}(\B,\B_\cc)$ 重合。一般说来，这大概期望过高，尽管明确的反而例子并不显然。但正如 \cite{GW} 所指出的，若 $M$ 紧，则有一个非常简单的论证给出 $\Hom_{A_K}(\B,\B_\cc)$ 与 $\Hom_{B_J}(\B,\B_\cc)$ 作为向量空间的同构，忽略它们的 Hermite 内积。论证很简单：若 $M$ 紧，则不必考虑某个扭转拓扑场论中的 $(\B,\B_\cc)$ 弦空间，而可以直接考虑 $(\B,\B_\cc)$ 系统的超对称基态。作为向量空间，这与任何可能的扭转超对称场论中的物理态空间相同；这一观察蕴含向量空间 $\Hom_{A_K}(\B,\B_\cc)\cong\Hom_{B_J}(\B,\B_\cc)$ 的等价。

然而，这一关系并\emph{不}延拓为 Hermite 内积之间的关系。没有理由期望 $\Hom_{A_K}(\B,\B_\cc)$ 的 Hermite 内积与式 \eqref{ngg} 所描述的 $\H_\gq$ 的内积一致。至于 $\Hom_{B_J}(\B,\B_\cc)$，要使它带有自然的 Hermite 结构，$Y$ 应具有一个在复结构 $J$ 中反全纯作用的 $\Z_2$ 对称。这样的对称的存在并不能由我们的假设推出，故一般说来 $\Hom_{B_J}(\B,\B_\cc)$ 没有 Hermite 结构。

在第 \ref{branecp} 节中，我们解释了：当 $M$ 具有由 $Y$ 的全纯极化 $\varPi$ 诱导的复结构 $\J$ 时，$\Hom_{A_K}(\B,\B_\cc)$ 作为向量空间与 $H^0_\J(M,K_M^{1/2}\otimes\frL)$ 相同。这一描述确实附带了一个描述 $\Hom_{A_K}(\B,\B_\cc)$ 的 Hermite 内积的配方，但该配方相当不显式，涉及一个一般难以计算的全纯三角形。人们并不期望这个配方一般会化为几何量子化所建议的通常定义 \eqref{ngg} 那样简单的东西。

于是我们有两种策略，把通过复化 $Y$ 对 $M$ 的膜量子化，与用复 Kähler 极化的量子化相比较。一种策略使用 $Y$ 的全纯极化，另一种使用 $Y$ 的超 Kähler 结构。

与其进一步讨论本文分析过的那类初等例子，我们将在更精细的例子的背景下讨论这两种途径；这个例子曾从超 Kähler 观点在 \cite{GW} 第 4 节中讨论过（那里给出了比我们这里所说的更多细节）。
